% Типовой расчёт № 1 по математическому анализу
% Лаборатория математики ФН1
% Компиляция: pdflatex tr-01-usloviya.tex

\documentclass[12pt, a4paper]{article}

\usepackage[T2A]{fontenc}
\usepackage[utf8]{inputenc}
\usepackage[russian]{babel}
\usepackage[left=25mm, right=20mm, top=20mm, bottom=20mm]{geometry}
\usepackage{amsmath, amssymb}
\usepackage{enumitem}

\pagestyle{plain}
\setlength{\parindent}{1.25cm}

\begin{document}

\begin{center}
  {\large\bfseries Типовой расчёт № 1}

  \vspace{0.3em}
  {\large Пределы и непрерывность}

  \vspace{0.5em}
  {\small Математический анализ, часть I · 1 курс, 1 семестр\\
  Лаборатория математики ФН1}
\end{center}

\vspace{1em}

\noindent
Номер варианта совпадает с порядковым номером студента в журнале группы.
Числовые данные вариантов — в файле \texttt{tr-01-varianty.csv}.

\section*{Задача 1. Предел последовательности}

Доказать по определению, что
\[
\lim_{n\to\infty}\frac{a n + 1}{n + b} = a,
\]
указав для $\varepsilon = 0{,}01$ конкретный номер $N(\varepsilon)$.

\section*{Задача 2. Пределы с неопределённостями}

Вычислить пределы, не используя правило Лопиталя:
\begin{enumerate}[label={\asbuk{enumi})}]
  \item $\displaystyle\lim_{n\to\infty}\left(\sqrt{n^2 + an} - n\right)$;
  \item $\displaystyle\lim_{n\to\infty}\left(\frac{n + c}{n + 1}\right)^{n}$;
  \item $\displaystyle\lim_{x\to 0}\frac{\sqrt{1 + kx} - 1}{x}$.
\end{enumerate}

\section*{Задача 3. Замечательные пределы}

Вычислить, сводя к первому и второму замечательным пределам:
\begin{enumerate}[label={\asbuk{enumi})}]
  \item $\displaystyle\lim_{x\to 0}\frac{\operatorname{tg} kx - \sin kx}{x^3}$;
  \item $\displaystyle\lim_{x\to\infty}\left(\frac{x + a}{x - a}\right)^{x}$.
\end{enumerate}

\section*{Задача 4. Сравнение бесконечно малых}

Определить порядок малости функции относительно $x$ при $x\to 0$:
\[
f(x) = \sqrt[3]{1 + x^2} - 1 + \sin^2 kx.
\]
Выделить главную часть вида $Cx^{\alpha}$.

\section*{Задача 5. Точки разрыва}

Найти точки разрыва функции, определить их тип и построить эскиз графика
в окрестности каждой точки:
\[
f(x) = \frac{x^2 - a^2}{x^2 - (a + 1)x + a}.
\]

\section*{Задача 6. Непрерывность с параметром}

При каком значении параметра $m$ функция непрерывна на всей числовой прямой?
\[
f(x) =
\begin{cases}
\dfrac{\sin mx}{x}, & x \ne 0, \\[2mm]
c, & x = 0.
\end{cases}
\]

\vspace{1.2em}

\section*{Требования к оформлению}

\begin{itemize}[nosep]
  \item условие переписывается полностью с числами своего варианта;
  \item в задаче~1 требуется именно доказательство по определению, а не вычисление предела;
  \item в задаче~5 эскизы графиков строятся на координатной сетке;
  \item ответ выделяется отдельной строкой.
\end{itemize}

\section*{Критерии оценивания}

\begin{center}
\begin{tabular}{|c|l|}
  \hline
  Баллы & За что начисляются \\
  \hline
  2 & задача решена верно, обоснование полное \\
  1 & верный метод, допущена вычислительная ошибка \\
  0 & метод выбран неверно либо решение отсутствует \\
  \hline
\end{tabular}
\end{center}

\vspace{0.5em}
\noindent
Расчёт зачтён при \textbf{8 баллах из 12} и успешной защите.

\vspace{0.7em}
\noindent\textbf{Сроки.} Выдача~--- 2-я учебная неделя,
сдача~--- \textbf{до конца 6-й недели}.

\end{document}
